\section{Conclusion and Future Work}
\label{sec:conclusion}

In this work, we presented \textbf{GranMerge}, a systematic empirical framework to deconstruct the \emph{Generalization-Granularity Trade-off} in multi-task adaptive model merging. 
By evaluating five nested levels of parameter resolution (Global, Layer-wise, Block-wise, Component-wise, and Tensor-wise) across four visual classification datasets using a 12-layer Vision Transformer, we mapped the precise performance boundaries of weight blending under test-time calibration constraints.

Our exhaustive sweeps revealed three fundamental findings. 
First, unconstrained test-time adaptation is highly susceptible to transductive overfitting, where increasing structural granularity (and thus parameter count) leads to progressive degradation of global multi-task generalization. 
Optimizing high-dimensional parameters on compact, unlabeled calibration streams ($N=256$) exploits local prediction entropy at the expense of the model's underlying representations. 
Second, optimizer choice fundamentally governs this overfitting: first-order Adam gradient descent exhibits rapid and severe generalization collapse, whereas black-box zero-order 1+1 Evolution Strategies maintain higher generalization, which we attribute to a combination of isotropic search boundaries and optimization sluggishness in high-dimensional spaces. 
Third, while simple soft L2 spatial-depth constraints (ESR and TV) can stabilize zero-order ES, they are insufficient to arrest the chaotic updates of gradient descent.

\textbf{Practical Guidelines for Practitioners.} 
For developers and researchers deploying adaptive model merging in the wild, our findings suggest several clear design guidelines:
\begin{enumerate}
    \item \textbf{Prefer low-dimensional parameters by default.} If unregularized adaptation is desired on small calibration streams, Global merging (Level 1) or uniform baselines represent the safest and most robust structural choices.
    \item \textbf{Use zero-order optimizers for high-dimensional search.} When adapting finer granularities on compact streams without explicit regularizations, zero-order optimizers (like 1+1 ES) offer a self-limiting optimization path that is naturally resilient to catastrophic overfitting.
    \item \textbf{Apply hard structural constraints for gradient descent.} If gradient-based methods must be used in high-dimensional spaces, simple soft L2 penalties are insufficient. Practitioners should employ hard structural boundaries (such as B-spline or low-degree polynomial parameterizations like PolyMerge) to restrict the optimization manifold.
\end{enumerate}

\textbf{Future Work and Strategic Directions.} 
We outline three highly promising avenues for future research to extend the GranMerge paradigm:
\begin{enumerate}
    \item \textbf{High-Fidelity Foundation Model Regimes.} A critical extension is evaluating GranMerge on large-scale, fully converged foundation models such as CLIP-Large or LLaMA-7B. In these high-fidelity regimes, the underlying representations are highly structured, clean, and stable. We hypothesize that the high-frequency parameter noise present in our low-fidelity, poorly converged experts is significantly mitigated in foundation models. Consequently, foundation models may exhibit a much more resilient generalization-granularity curve, potentially allowing fine-grained adaptive merging to outperform uniform baselines by precisely aligning well-defined feature spaces without collapsing under transductive updates.
    \item \textbf{Exploring Semantically Richer Surrogate Losses.} While prediction entropy is a convenient differentiable objective for test-time adaptation, our analysis of surrogate loss misalignment (Section~\ref{sec:experiments}) shows its severe limitations. Future work should investigate semantically richer surrogate objectives that capture representation-level structures rather than simple marginal output confidence. Concrete starting points include: (a) \emph{self-supervised contrastive objectives} that preserve representation consistency across augmented views of the calibration stream, (b) \emph{multi-modal joint calibration losses} that leverage cross-modal alignment (e.g., image-text similarity) to anchor weight blending, and (c) \emph{soft feature-alignment objectives} that minimize CKA (Centered Kernel Alignment) drift from base model representations.
    \item \textbf{Hybrid Optimization Paradigms.} Finally, future work should explore hybrid optimizers that combine the strengths of both optimizer families: leveraging the rapid, coordinate-aligned convergence of first-order gradient descent in low-dimensional subspaces (such as L1/L2) while deploying the robust, self-limiting isotropic walks of zero-order search (such as 1+1 ES or Guided ES) in high-dimensional subspaces to prevent catastrophic representation collapse.
\end{enumerate}
We hope our systematic study provides a solid empirical foundation for more robust, deployment-ready adaptive model merging frameworks.
